\section{Conclusion}
Reference-derived phase plans connect a deployment's numerical behavior to its parallel execution. They combine the producer history, physical KV range, type-specific arithmetic and ownership of logical outputs and state. Whole-graph and local-attention implementations place the cuts differently while preserving the reference-sensitive computation within each group.

CUDA and Vulkan studies support this approach across model structures, operating environments and quantized paths, including mixed-Q4 at 32B. Controlled comparisons separate packing's contribution from geometry-dependent local cuts. Complete requests then determine the useful deployment: speculative execution improves the three Linux service panels, grouped ordinary generation is preferable in the tested 32B configuration, and shorter proposals remain preferable in the live width study despite the aggregate equal-work advantage of eight-query execution. Phase plans make these choices available within an established numerical reference.
